\scanpage{4}
% archive foliation: 3
\uline{Lemme 2} Soit $L_*$ un $k$-module $\frac12$ \unsure{simplicial},
\inserted{(\unsure{projectif})} \deleted{flasque}, \inserted{\unsure{$q\geq q^{\prime}\geq0$}} \deleted{\unsure{$q>0$}}, $M$ un $k$-module, \unsure{alors}
\[
\unsure{\operatorname{Hom}_{\widehat\Delta_k}}
\bigl(L_*\otimes_k\unsure{\Lambda^q\Psi_{* k}},
\unsure{\Lambda^{q'}\Psi_{* k}}\otimes_k M\bigr)=0
\]
\inserted{(on utilise cela \unsure{pour un complexe nul en degrés} $<q$)}.

\deleted{\unsure{Soit} \illegible{7}}
\inserted{le complexe de chaînes associé à}
\deleted{\illegible{4}}
$L_*\otimes_k\unsure{\Lambda^q\Psi_{* k}}$ \unsure{est}
\deleted{\illegible{3}}
\deleted{\unsure{un} \illegible{5}} \unsure{de la forme}
\[
\cdots\longrightarrow L'_{q+1}\longrightarrow L'_q\longrightarrow0\longrightarrow\cdots
\]
tandis que $\unsure{\Lambda^{q'}\Psi_{* k}}\otimes_k M$ est $M$ placé au \unsure{degré}
$q'$ \inserted{\unsure{i.e. $M[q']$}}. \unsure{Si} $q'=q$, il \unsure{s'agit} de \illegible{1}
$\unsure{\operatorname{Hom}}(L'_*,M\unsure{[q]})$
$=0$ \unsure{pour un morphisme de} \illegible{2} \unsure{de la forme}
\diagram{\begin{tikzcd}
\unsure{L'_{q+1}}\arrow[r,"\unsure{d}"]\arrow[d]&\unsure{L'_q}\arrow[r]\arrow[d]&0\\
0\arrow[r]&M\arrow[r]&0
\end{tikzcd}}
alors $\pi_{\bullet}(L_{\bullet})=0$ $\Rightarrow$ $L'_{q+1}\xrightarrow{d}L'_q$ est surjectif,
donc $L'_q\to M$ est nul, \unsure{donc le hom. de}
\unsure{complexes} \inserted{de chaînes} $\unsure{L'_*\to M[q]}$, cqfd.

\uline{Théorème} Soit $n\geq0$. Il existe un \unsure{homomorphisme}
de complexes de \unsure{modules dans} $\unsure{\widehat\Delta_k}$
\diagram{\begin{tikzcd}[column sep=small]
0\arrow[r]&\unsure{\Gamma^n\Phi_*}\arrow[r]\arrow[d,"u_0"]
&\unsure{\Gamma^{n-1}\Phi_*\otimes\Psi_*}\arrow[r]\arrow[d,"u_1"]
&\unsure{\Gamma^{n-2}\Phi_*\otimes\Lambda^2\Psi_*}\arrow[r]\arrow[d,"u_2"]
&\cdots\arrow[r]&\unsure{\Lambda^n\Psi_*}\arrow[r]\arrow[d,"u_n"]&0\arrow[d]\\
0\arrow[r]&\unsure{\Phi_*}\arrow[r,"\unsure{\wedge T}"']
&\unsure{\Lambda^2\Omega_*}\arrow[r,"\unsure{\wedge T}"']
&\unsure{\Lambda^3\Omega_*}\arrow[r]&\cdots\arrow[r,"\unsure{\wedge T}"']
&\unsure{\Lambda^{n+1}\Omega_*}\arrow[r,"\unsure{\wedge T}"']&\unsure{\Lambda^{n+2}\Omega_*}
\end{tikzcd}}
(\unsure{nb}: on écrit $\unsure{\Phi_*}$ et $\unsure{\Psi_*}$ pour
$\unsure{\Phi_{* k}}$, $\unsure{\Psi_{* k}}$, et $\unsure{\Omega}$
\unsure{l'algèbre} $\unsure{\Omega_k}$) tel que \unsure{$u_n$ soit l'isom. de}
$\unsure{\Lambda^n\Psi_*}$ sur
$\unsure{\operatorname{Ker}(\Lambda^{n+1}\Omega_*\to\Lambda^{n+2}\Omega_*)
=d\,Z(\Lambda^n\Omega_*)}$
